\documentclass[12pt,a4paper,oneside]{ctexart}
\usepackage{amsmath,amsthm,amssymb,bm,graphicx,hyperref,mathrsfs}

\title{\textbf{作业}}
% \author{黄伟宁}
\date{\today}
\linespread{1.5}
\newcounter{problemname}
\newenvironment{problem}{\stepcounter{problemname}\par\noindent\textbf{题目\arabic{problemname}. }}{\\\par}
\newenvironment{solution}{\par\noindent\textbf{解答. }}{\\\par}
\newenvironment{note}{\par\noindent\textbf{题目\arabic{problemname}的注记. }}{\\\par}
\usepackage{geometry}
\geometry{left=2.54cm,right=2.54cm,top=3.18cm,bottom=3.18cm}

\begin{document}

\maketitle
% \begin{problem}
% 这里是题目.
% \end{problem}

% \begin{solution}
%     这里是解答.
% \end{solution}

% \begin{note}
%     这里是注记.
% \end{note}

% \begin{problem}
% \end{problem}
\begin{problem}
    已知函数$ f(x)=\log_a(\frac{1}{2}ax^2-x+\frac{1}{2}) $ 在区间$ [2,3] $ 上恒为正则实数$ a $ 的取值范围是$ \underline{\qquad\qquad} $ .
\end{problem}

\begin{problem}
    已知函数$ f(x)=\begin{cases}
        |\log_2 x|,0<x\leqslant 4,\\
        \frac{2}{3}x^2-8x+\frac{70}{3},x>4,
    \end{cases} $若$ a,b,c,d $ 互不相同,且$ f(a)=f(b)=f(c)=f(d) $ ,则$ abcd $ 的取值范围是$ \underline{\qquad\qquad} $ .
\end{problem}

\begin{problem}
    若$ \log_4(x+2y)+\log_4(x-2y)=1 $,则$ |x|-|y| $ 的最小值是$ \underline{\qquad\qquad} $ . 
\end{problem}

\begin{problem}
    已知函数$ y=f(x) $ 的图像既关于点$ (1,1) $ 中心对称,又关于直线$ x+y=0 $ 轴对称.若当$ x\in (0,1) $ 时,$ f(x)=log_2(x+1) $,则$ f(log_2 10) $ 的值为$ \underline{\qquad\qquad} $ .
\end{problem}

\begin{problem}
    设函数$ f(x)=|\lg(x+1)| $ ,实数$ a,b(a<b) $ 满足$ f(a)=f(-\frac{b+1}{b+2}),f(10a+6b+21)=4\lg 2 $,求$ a,b $ 的值. 
\end{problem}

\begin{problem}
    设$ f(x) $ 是定义在区间$ (0,+\infty) $ 上的单调函数,对于任意的$ x>0 $ ,均有$ f(x)>-\frac{4}{x},f(f(x)+\frac{4}{x})=3, $ 则$ f(8)=\underline{\qquad\qquad}. $ 
\end{problem}

\begin{problem}
    设函数$ f(x) $ 对一切$ x>0 $ 有定义,且满足:

    1. $ f(x) $ 在区间$ (0,+\infty) $ 上严格单调递增;

    2. 对一切$ x>0 $ ,有$ f(x)>-\frac{1}{x} $;

    3. 对一切$ x>0 $ ,有$ f(x)f(f(x)+\frac{1}{x})=1 $ .

    则$ f(1)=\underline{\qquad\qquad} $ .
\end{problem}

\begin{problem}
    已知$ a<b $ ,二次函数$ y=ax^2+bx+c\geqslant 0 $ 对任意实数$ x $ 恒成立,则$ M=\frac{a+2b+4c}{b-a} $ 的最小值是$ \underline{\qquad\qquad} $ .
\end{problem}

\begin{problem}
    已知奇函数$ f(x) $ 在定义域$ [-3,3] $ 上是减函数,且$ f(x^2-2x)+f(x-2)<0 $,求实数$ x $的取值范围.
\end{problem}

\begin{problem}
    设$ f(x)=x^2+bx+c(b,c\in \mathbb{R}) .$当$ |x|\geqslant 2 $ 时,$ f(x)\geqslant 0, $且$ f(x) $在区间$ (2,3] $上的最大值为$ 1 $ ,求$ b^2+c^2 $ 的最大值和最小值.
\end{problem}

\begin{problem}
    已知函数$ f(x) $ 的定义域为$ \mathbb{R} $ ,满足:1. $ f(x) $ 是偶函数;2. 对于任意的$ x\in \mathbb{R} ,$都有$ f(x+4)=f(x), $且当$ x\in [0,2] $时,$ f(x)=x+2 $.则直线$ y=4 $ 与函数$ f(x) $ 的图像交点中最近的两点之间的距离为$ \underline{\qquad\qquad} $.
\end{problem}

\begin{problem}
    设函数$ f(x) $ 定义在$ \mathbb{R} $ 上,给出下述三个命题:

    (1) 满足条件$ f(x+2)+f(2-x)=4 $ 的函数图象关于点$ (2,2) $ 对称;

    (2) 满足条件$ f(x+2)=f(2-x) $ 的函数图象关于直线$ x=2 $ 对称;3

    (3) 函数$ f(x-2) $ 与$ f(-x+2) $ 在同一坐标系中的图象关于直线$ x=2 $ 对称.

    其中,真命题的个数是$ \underline{\qquad\qquad} $ .
\end{problem}

\begin{problem}
    设$ f(x) $ 是定义在$ \mathbb{R} $ 上以$ 2 $ 为周期的偶函数,在区间$ [0,1] $ 上严格递减,且满足$ f(\pi)=1,f(2\pi)=2, $ 则不等式组$ \begin{cases}
        1\leqslant x\leqslant 2,\\
        1\leqslant f(x)\leqslant 2
    \end{cases} $的解集为$ \underline{\qquad\qquad} $.
\end{problem}

\begin{problem}
    定义在$ \mathbb{R} $ 上的函数$ f(x) $ 的图象关于点$ (-\frac{3}{4},0) $ 中心对称,对任意的实数$ x $,都有$ f(x)=-f(x+\frac{3}{2}) ,$且$ f(-1)=1,f(0)=-2, $求$ f(1)+f(2)+f(3)+\cdots+f(2009) $的值. 
\end{problem}

\begin{problem}
    已知二次函数$ y=f(x)=x^2+bx+c $ 的图象过点$ (1,13) $,且函数$ y=f(x-\frac{1}{2}) $是偶函数.

    (1)求$ f(x) $ 的解析式.

    (2)函数$ y=f(x) $ 的图象上是否存在这样的点,其横坐标是正整数,纵坐标是一个完全平方数? 如果存在,求出这样的点的坐标;如果不存在,请说明理由.
\end{problem}

\end{document}